% TOP_PROBLEM: {"id":30007525,"problem_number":"LOCAL-30007525","title":"Heterogeneous monetary transmission","category_id":21,"category":{"id":21,"name":"economics","display_name":"Economics","description":"Open economic theory statements, published research agendas, and proposed scoped specifications; question provenance and historical priority are qualified per record.","slug":"economics","order_index":21,"created_at":"2026-10-08T00:00:00Z"},"status":"research_question","external_url":null,"legacy_ids":[],"published":false,"statement":"Identify consumption responses to a monetary shock across predetermined household balance-sheet exposures and assess whether they aggregate to observed macroeconomic responses.","economics":{"original_id":14,"field":"Inflation and monetary policy","kind":"identification","formalization_basis":"adapted_research_question","source_status":"research_agenda","priority_status":"earliest_found","priority_note":"This is the earliest explicit relevant statement verified in this audit, not a claim of original priority; older literature and earlier versions may contain antecedents.","earliest_verified_reference":{"authors":["Janet L. Yellen"],"title":"Macroeconomic Research After the Crisis","year":2016,"url":"https://www.federalreserve.gov/newsevents/speech/yellen20161014a.htm","doi":null,"locator":"Section “Heterogeneity,” paragraphs on aggregate demand and financial constraints","evidence_summary":"The speech explicitly asks how individual differences affect aggregate demand and says heterogeneous financial constraints and firm linkages remain insufficiently understood empirically and theoretically."},"current_reference":{"authors":["Adrien Auclert","Matthew Rognlie","Ludwig Straub"],"title":"Fiscal and Monetary Policy with Heterogeneous Agents","year":2025,"url":"https://web.stanford.edu/~aauclert/annual_review_hank.pdf","doi":"10.1146/annurev-economics-091624-044646","locator":"“Nominal assets,” p. 16 of the April 2025 author manuscript","evidence_summary":"The review explicitly says the empirical magnitude of the Fisher redistribution channel remains open to debate; it does not claim that all channels of monetary transmission are unknown."},"formalization_note":"The household target narrows the wider household/firm/bank catalogue question. The current review explicitly flags the magnitude of the nominal redistribution channel as debated."},"statusQualification":"Published question or research agenda verified at the cited locator. The mathematical specification is adapted for this catalogue and includes additional assumptions; source publication does not establish first-ever priority or certify this exact model remains unsolved. The household target narrows the wider household/firm/bank catalogue question. The current review explicitly flags the magnitude of the nominal redistribution channel as debated.","statusReviewedAt":"2026-10-08"}
% FIRST_POSTED: 2026-10-08
% ENTRY_KIND: statement_only
% !TeX program = lualatex
\documentclass[11pt]{article}
\usepackage{fontspec}
\usepackage[a4paper,margin=25mm]{geometry}
\usepackage{amsmath,amssymb,mathtools}
\usepackage{xurl}
\usepackage[hidelinks]{hyperref}
\newcommand{\sref}[2]{\hyperlink{source:#1:#2}{[S#2]}}
\setlength{\emergencystretch}{3em}
\begin{document}
% problemId:30007525
\section*{Heterogeneous monetary transmission}\label{30007525}
\subsection{Definitions and mathematical statement}
Fix households with predetermined type \(h_i=(a_i^{L},a_i^{I},d_i,n_i)\), where liquid assets, illiquid assets, debt, and nominal asset exposure are measured before a monetary intervention. Let \(m_t\) be an externally identified unexpected increase in the policy rate, and \(c_{i,t+H}(m)\) household real consumption under shock size \(m\), including equilibrium income and price adjustments. Define the conditional response
\[
\tau_H(h)=\left.\frac{\partial E[c_{i,t+H}(m)\mid h_i=h]}{\partial m}\right|_{m=0}.
\]
Identify \(\tau_H(h)\) and its aggregation \(\int\tau_H(h)\,dF(h)\), where \(F\) is the population distribution of predetermined types. Provide instruments that isolate monetary surprises from central-bank information about fundamentals, and address nonrandom refinancing, unobserved insurance, and changes in household composition. Separate direct interest-rate and nominal-revaluation channels from induced labor-income and asset-price responses through a stated structural model or independently credible channel interventions. Test whether identified heterogeneous effects aggregate to observed consumption responses across episodes. The target includes uncertainty about the quantitative Fisher redistribution channel; it does not assume that all transmission mechanisms are unknown or that conditioning on wealth alone identifies a causal channel.

\par\textbf{Formulation and scope.} The household target narrows the wider household/firm/bank catalogue question. The current review explicitly flags the magnitude of the nominal redistribution channel as debated.

\par\textbf{Open-question provenance.} An explicit question or research agenda was verified at the source locator. This does not by itself certify that every later version remains unresolved \sref{30007525}{1}.

\par\textbf{Historical priority.} This is the earliest explicit relevant statement verified in this audit, not a claim of original priority; older literature and earlier versions may contain antecedents.

\subsection{Short English statement}
Identify consumption responses to a monetary shock across predetermined household balance-sheet exposures and assess whether they aggregate to observed macroeconomic responses.

\subsection{Sources}
\begin{itemize}\raggedright
\item[\textnormal{[S1]}]\hypertarget{source:30007525:1}{} Janet L. Yellen. 2016. Macroeconomic Research After the Crisis. Section “Heterogeneity,” paragraphs on aggregate demand and financial constraints.
\url{https://www.federalreserve.gov/newsevents/speech/yellen20161014a.htm}.
{\small\textit{Earliest verified question/agenda reference; historical priority qualified below. The speech explicitly asks how individual differences affect aggregate demand and says heterogeneous financial constraints and firm linkages remain insufficiently understood empirically and theoretically.}}

\item[\textnormal{[S2]}]\hypertarget{source:30007525:2}{} Adrien Auclert, Matthew Rognlie, Ludwig Straub. 2025. Fiscal and Monetary Policy with Heterogeneous Agents. “Nominal assets,” p. 16 of the April 2025 author manuscript.
\url{https://web.stanford.edu/~aauclert/annual_review_hank.pdf}.
DOI: 10.1146/annurev-economics-091624-044646.
{\small\textit{Supporting or current literature; see evidence qualification. The review explicitly says the empirical magnitude of the Fisher redistribution channel remains open to debate; it does not claim that all channels of monetary transmission are unknown.}}
\end{itemize}
\end{document}
